\documentclass{article}
\usepackage[a4paper, total={6.5in, 10in}]{geometry}
\setlength{\parindent}{0pt}
\usepackage{tcolorbox}
\tcbset{colback=white!94!black, colframe=black!60!white, coltitle = white, boxrule=0.8pt, arc=2mm}
\newtcolorbox{boxs}[2][]{title= #2,#1}
\usepackage{datetime2}
\usepackage[british]{babel}
\usepackage{amsfonts}
\usepackage{amsmath}

\title{\textbf{Real Analysis HW 4}}
\author{Robert O'Connell}
\date{\today}
\begin{document}
\maketitle
\vspace{5mm}

\subsection*{BS - Page 67 Exercise 13}
Use the Squeeze Theorem to determine the limits of 
\begin{itemize}
    \item \(\left(n^{\frac{1}{n^2}}\right)\)
    \item \(\left( (n!)^{\frac{1}{n^2}}\right)\)
\end{itemize}

\subsubsection*{(a)}
\[
n^{\frac{1}{n^2}} = e^{\ln(n^{\frac{1}{n^2}})} = e^{\frac{\ln(n)}{n^2}}
\]
Now we only need to condsider \(\frac{\ln(n)}{n^2}\)
\\

Since \(\ln(n) \geq 0 \ \forall \ n\) and \(\ln(n) \leq n \ \forall \ n\), proofs at the end of this question, we have
\[
0 \leq \frac{\ln(n)}{n^2} \leq \frac{n}{n^2} = \frac{1}{n}
\]

Then by Squeeze Theorem, since \(\frac{1}{n} \rightarrow 0\), we have that \(\frac{\ln(n)}{n^2}\) converges to \(0\)
\\

Thus \(e^{\frac{\ln(n)}{n^2}} = n^{\frac{1}{n^2}}\) converges to \(e^0 = 1\)
\\

\subsubsection*{(b)}

Similarly for the second limit we get,
\[
(n!)^\frac{1}{n^2} = e^{\frac{\ln(n!)}{n^2}}
\]

\[
\frac{\ln(n!)}{n^2} = \frac{1}{n^2} (\ln(n) + \ln(n-1) + \dots + \ln(1))
\]

Then since \(\ln(n) \leq \ln(n+1) \ \forall \ n\) we get,
\[
0 \leq \frac{\ln(n!)}{n^2} \leq \frac{1}{n^2}(n\ln(n))
\]
\[
0 \leq \frac{\ln(n!)}{n^2} \leq \frac{1}{n}(\ln(n)) 
\]

We have \(\ln(n) \leq n \Rightarrow \ln(\sqrt{n}) \leq \sqrt{n} \Rightarrow \ln(n) \leq 2\sqrt{n}\)
\[
0 \leq \frac{\ln(n!)}{n^2} \leq \frac{2\sqrt{n}}{n} = \frac{2}{\sqrt{n}}
\]
As \(n \rightarrow \infty, \sqrt{n} \rightarrow \infty\), since we know \(\frac{1}{n} \rightarrow 0\) 
then \(\frac{2}{\sqrt{n}} \rightarrow 0\)

Thus \(\frac{\ln(n!)}{n^2}\) too converges to \(0\), by Squeeze Theorem
\\

Thus \(e^{\frac{\ln(n!)}{n^2}} = (n!)^{\frac{1}{n^2}}\) converges to \(e^0 = 1\)
\\

\subsubsection*{Proofs of properties used}

\begin{itemize}
    \item I will prove \(\ln(n) \geq 0 \ \forall \ n\) by induction,

    Base Case:
    \[
    \ln(1) = 0 \geq 0
    \]

    Induction Step:

    Assume it holds for \(n = k\)
    \begin{align*}
        \ln(k) \geq 0 \\
        k \geq e^0 = 1 \\
        k + 1 \geq 2 \\
        \ln(k+1) \geq \ln(2) \geq 0    
    \end{align*}
    Thus \(\ln(n) \geq 0 \ \forall \ n\) 
    \\

    \item \(\ln(n) \leq n \ \forall \ n\) will also be proven by induction

    Base Case:
    \[
    \ln(1) = 0 \leq 1
    \]

    Induction Step:

    Assume it holds for \(n=k\)
    \begin{align*}
        \ln(k) \leq k \\
        k \leq e^k \\
        ke \leq e^{k+1} \\
    \end{align*}

    Also \(k + 1 \leq ke \ \forall \ k\)
    
    \begin{itemize}
        \item Base Case:
    \[
    2 \leq e
    \]

    Induction Step:

    Assume it holds for \(k = m\)
    \begin{align*}
        m + 1 \leq me \\
        m + 1 + e \leq e(m+1) \\
        m+ 2 \leq m+ 1 + e \leq e(m+1)
    \end{align*}
    \end{itemize}

    Then 
    \begin{align*}
        k + 1 \leq ke \leq e^{k +1 }\\
        \ln(k+1) \leq k + 1
    \end{align*}
    Thus by induction, \(\ln(n) \leq n \ \forall \ n\)
\end{itemize}









\subsection*{BS - Page 80 Exercise 5}
Let \(X = \{x_n\}\) and \(Y = \{y_n\}\) be given sequences and let the "shuffled" 
sequence \(Z= \{z_n\}\) be defined by \(z_1 = x_1, z_2 = y_1, \dots, z_{2n-1} = x_n, z_{2n} = y_n, \dots\) 

Show that \(Z\) is convergent iff both \(X\) and \(Y\) are convergent and \(\lim X = \lim Y\)
\\

(\(\Leftarrow\))

Assume \(\{x_n\}, \{y_n\}\) convergent, both with limit \(L\)
\[
\Rightarrow \forall \ \epsilon >0 \ \exists \ N_1, N_2 \in \mathbb{N} \text{ such that }
\]
\begin{itemize}
    \item \(|x_n - L| < \epsilon  \ \forall \ n \geq N_1\)
    \item \(|y_n - L| < \epsilon  \ \forall \ n \geq N_2\)
\end{itemize}

\(\{z_n\}\) convergent  with limit \(z\) if \(\forall \ \epsilon > 0 \ \exists \ N_3 \in \mathbb{N} 
\text{ such that } |z_n -z| < \epsilon \ \forall \ n \geq N_3 \)

There are two cases for each \(z_n\)
\[
z_n = \begin{cases}
    x_{\frac{n+1}{2}}, \quad n \text{ odd} \\
    y_{\frac{n}{2}}, \quad n \text{ even}
\end{cases}
\]

\(\forall \ n \geq \max\{N_1, N_2\}\), we have \(|x_n - L| < \epsilon \) and \(|y_n -L| < \epsilon \)

\(\Rightarrow \forall \ n \geq \max\{2N_1, 2N_2\}\), we have \(|x_{\frac{n+1}{2}} - L| < \epsilon \) if \(n\) is odd
and \(|y_{\frac{n}{2}} -L| < \epsilon \) if \(n\) is even.

\[
\Rightarrow \forall \ \epsilon >0 \ \exists \ N_3 \in \mathbb{N} \text{ such that } |z_n - L| < \epsilon  \ \forall \ n \geq N_3
\]

\(\Rightarrow \{z_n\}\) convergent with limit \(L\)
\\

(\(\Rightarrow\))

Assume \(\{z_n\}\) convergent with limit \(L\)
\[
\forall \ \epsilon >0 \ \exists \ N_1 \in \mathbb{N} \text{ such that } |z_n - L| < \epsilon  \ \forall \ n \geq N_1
\]

There are two cases for each \(z_n\)
\[
z_n = \begin{cases}
    x_{\frac{n+1}{2}}, \quad n \text{ odd} \\
    y_{\frac{n}{2}}, \quad n \text{ even}
\end{cases}
\]


\(\forall \ n \geq N_1,\) of the form \(n = 2k + 1, k \in \mathbb{N}\) \(\Rightarrow |x_{k+1} - L| < \epsilon \) 
\(\forall \ k \geq \frac{N_1 - 1}{2}\)
\\

\(\forall \ n \geq N_1,\) of the form \(n = 2k, k \in \mathbb{N}\) \(\Rightarrow |y_k - L| < \epsilon \) 
\(\forall \ k \geq \frac{N_1}{2}\)


\[
\Rightarrow \forall \ \epsilon >0 \ \exists \ N_2 \in \mathbb{N} \text{ such that }
\]
\begin{itemize}
    \item \(|x_n - L| < \epsilon  \ \forall \ n \geq N_2\)
    \item \(|y_n - L| < \epsilon  \ \forall \ n \geq N_2\)
\end{itemize}


\subsection*{BS - Page 80 Exercise 7}
Establish the convergence and find the limits of the following sequences
\begin{itemize}
    \item \(\left(\left(1+ \frac{1}{n^2}\right)^{n^2}\right)\)
    \item \(\left(\left(1+ \frac{1}{2n}\right)^{n}\right)\)
    \item \(\left(\left(1+ \frac{1}{n^2}\right)^{2n^2}\right)\)
    \item \(\left(\left(1+ \frac{2}{n}\right)^{n}\right)\)
\end{itemize}

We know \((1+\frac{1}{k})^k\) converges to \(e\)
\begin{itemize}
    \item We have \(\{x_n\} = (1+\frac{1}{n^2})^{n^2}\)
    \\
    Let \(\{x_k\} = (1+ \frac{1}{k})^k\) \\
    Then \(\{x_n\}\) is a subsequence of \(\{x_k\}\), with the indices being the perfect squares,
    \\ 
    Since we know \(\{x_k\}\) converges to \(e\), we also have that it's subsequence \(\{x_n\}\) converges to \(e\)

    \item We have \(\{x_n\} = (1+ \frac{1}{2n})^n = ((1+ \frac{1}{2n})^{2n})^{\frac{1}{2}}\)
    \\
    Let \(\{x_m\} = (1+ \frac{1}{2n})^{2n}\)
    \\
    Similarly, we have \(\{x_m\}\) as a subsequence of \(\{x_k\}\), with the indices being
    multiples of \(2\), thus \(\{x_m\}\) converges to \(e\)
    \\
    Since \(\{x_m\}\) converges to \(e\), we have that \(\{x_n\} = (\{x_m\})^{\frac{1}{2}}\) converges to \(\sqrt{e}\)

    \item We have \(\{x_n\} = (1+ \frac{1}{n^2})^{2n^2} = ((1+ \frac{1}{n^2})^{n^2})^{2}\)
    \\
    Let \(\{x_m\} = (1+ \frac{1}{n^2})^{n^2}\)
    \\
    From earlier, we have \(\{x_m\}\) converges to \(e\)
    \\
    Since \(\{x_m\}\) converges to \(e\), we have that \(\{x_n\} = (\{x_m\})^2\) converges to \(e^2\)

    \item We have \(\{x_n\} = (1+ \frac{2}{n})^n = ((1+ \frac{1}{\frac{n}{2}})^{\frac{n}{2}})^2\)
    \\
    Let \(\{x_m\} = (1+ \frac{1}{\frac{n}{2}})^{\frac{n}{2}}\)
    \\
    We have \(\{x_m\}\) converges to \(e\), as when \(k \rightarrow \infty\) in \(\{x_k\}\), 
    \(m \rightarrow \infty\) in \(\{x_m\}\)
    \\
    Since \(\{x_m\}\) converges to \(e\), we have that \(\{x_n\} = (\{x_m\})^{2}\) converges to \(e^2\)
\end{itemize}

\subsection*{BS - Page 80 Exercise 11}
Suppose that \(x_n \geq 0 \) for all \(n \in \mathbb{N}\) and that \(\lim((-1)^n x_n)\)
exists. Show that \(\{x_n\}\) converges
\\

Let \(\lim((-1)^n x_n) = x\)
\[
\Rightarrow \forall \ \epsilon > 0 \ \exists \ N \in \mathbb{N} \text{ such that } |(-1)^nx_n - x| < \epsilon \ \forall \ n \geq N
\]

Then for any even \(k\), \(k\geq N\)
\[
\Rightarrow |(-1)^k x_k - x| < \epsilon 
\]
\[
\Rightarrow |x_k - x| < \epsilon 
\]


Also for any odd \(p\), \(p \geq N\)
\[
\Rightarrow |(-1)^px_p - x| < \epsilon 
\]
\[
\Rightarrow |x_p + x| < \epsilon 
\]

Since \(x_n \geq 0 \ \forall \ n\) and \(\epsilon  > 0\)
\[
\Rightarrow x \geq 0, \quad x \leq 0
\]
Because in the first equation if \(x < 0\) it implies \(x_k - x > 0\) meaning you could pick \(\epsilon  = \frac{-x}{2}\)
to get a contradiction, and in the second if \(x > 0\) implies \(x_p + x > 0\) meaning you could pick \(\epsilon  = \frac{x}{2}\)
to get a contradiction

Thus \(x=0\)
\\

Then
\[
\Rightarrow \forall \ \epsilon > 0 \ \exists \ N \in \mathbb{N} \text{ such that } |(-1)^nx_n | < \epsilon \ \forall \ n \geq N
\]
\[
\Rightarrow x_n < \epsilon 
\]
Hence \(\{x_n\}\) converges to \(0\)
\end{document}